% Autorska priprava k Realisticky.cz 020411; ulohy jsou autorske a obmenene.
% Zamerne nenahrazuje vyklad presnymi kopiemi cizich textu/prikladu.
\documentclass[11pt,a4paper]{article}
\input{../../sablony/pripravy/priprava.sty}
\usepackage{tikz}
\nastavPripravu{1. ročník}{Rovnice a nerovnice}{89}{Jednoduché nerovnice s absolutní hodnotou}
\begin{document}
\blok{Návaznost a motivace}
Absolutní hodnotu chápeme jako vzdálenost. Výrok $|x-a|<r$ znamená, že $x$ leží blíže než $r$ k číslu $a$; opačná nerovnost znamená vnější část číselné osy.
\blok{Princip a odvození}
Pro $r>0$: $|A(x)|<r$ právě tehdy, když $-r<A(x)<r$ (průnik podmínek), kdežto $|A(x)|>r$ vede na $A(x)<-r$ \emph{nebo} $A(x)>r$ (sjednocení). Podobně pro neostré nerovnosti. Ukázky zde necháváme jednoduché, bez dvou modulů.
\blok{Řešené příklady: od jednoduchého ke složitějšímu}
\begin{prikladReseniBox}[title={Příklad 1}]
Vyřešte $|x-2|<3$.
\tcblower
\begin{align*}
-3&<x-2<3 \\
 -1&<x<5 \\
 K&=(-1;5)
\end{align*}
\end{prikladReseniBox}
\begin{prikladReseniBox}[title={Příklad 2}]
Vyřešte $|2x+1|\ge5$.
\tcblower
\begin{align*}
2x+1&\le-5\quad\text{nebo}\quad 2x+1\ge5 \\
 x&\le-3\quad\text{nebo}\quad x\ge2 \\
 K&=(-\infty;-3\rangle\cup\langle2;+\infty)
\end{align*}
\end{prikladReseniBox}
\poznamka{Použij číselnou osu a význam vzdálenosti. Složitější nerovnice se dvěma moduly jsou jen nadstavba, nikoliv cíl této hodiny.}
\blok{Samostatné procvičování}
\begin{zadaniReseniBox}[title={Úloha 1}]
$|x|<4$.
\tcblower
\begin{align*}
-4&<x<4 \\
 K&=(-4;4)
\end{align*}
\end{zadaniReseniBox}
\begin{zadaniReseniBox}[title={Úloha 2}]
$|x-1|\ge2$.
\tcblower
\begin{align*}
x-1&\le-2\quad\text{nebo}\quad x-1\ge2 \\
 x&\le-1\quad\text{nebo}\quad x\ge3
\end{align*}
\end{zadaniReseniBox}
\begin{zadaniReseniBox}[title={Úloha 3}]
$|x+4|\le1$.
\tcblower
\begin{align*}
-1&\le x+4\le1 \\
 -5&\le x\le-3 \\
 K&=\langle-5;-3\rangle
\end{align*}
\end{zadaniReseniBox}
\blok{Volitelná rezerva}
\begin{zadaniReseniBox}[title={Úloha 1}]
Vyřešte $|x+2|\le0$.
\tcblower
\begin{align*}
|x+2|&\ge0 \\
 x+2&=0 \\
 K&=\{-2\}
\end{align*}

Při ostré nerovnosti $|x+2|<0$ by naopak bylo $K=\varnothing$.
\end{zadaniReseniBox}
\blok{Domácí úkol}
\begin{zadaniReseniBox}[title={Domácí příprava}]
\begin{enumerate}[label=\alph*)]
\item
$|3x|<6$.
\item
$|x-2|>4$.
\item
$|x+3|\le3$.
\item
$|x-4|\le2$.
\item
$|2x+1|>5$.
\item
$|x+1|\ge0$.
\end{enumerate}
\tcblower
\begin{enumerate}[label=\alph*)]
\item
\begin{align*}
-6&<3x<6 \\
 -2&<x<2 \\
 K&=(-2;2)
\end{align*}
\item
\begin{align*}
x-2&<-4\quad\text{nebo}\quad x-2>4 \\
 x&<-2\quad\text{nebo}\quad x>6
\end{align*}
\item
\begin{align*}
-3&\le x+3\le3 \\
 -6&\le x\le0 \\
 K&=\langle-6;0\rangle
\end{align*}
\item
$\begin{aligned}
-2&\le x-4\le2\\2&\le x\le6\\K&=\langle2;6\rangle
\end{aligned}$
\item
$\begin{aligned}
2x+1&<-5\quad\text{nebo}\quad2x+1>5\\x&<-3\quad\text{nebo}\quad x>2\\K&=(-\infty;-3)\cup(2;+\infty)
\end{aligned}$
\item
$\begin{aligned}
|x+1|&\ge0\quad\text{pro každé }x\in\mathbb R\\K&=\mathbb R
\end{aligned}$
\end{enumerate}
\end{zadaniReseniBox}
\poznamka{Orientační časové rozvržení: 5 + 12 + 12 + 10 + 6 = 45 min. Přednost má pochopení jádra; rezerva není součástí časového součtu. Domácí úlohy jsou určeny na samostatnou práci, řešení zůstávají učiteli.}
\end{document}
